\writeSectionFrame{\presentationTitle}{Problem bei der Serialisierung}
\begin{frame}[fragile]{Serialisierung}
	\begin{exampleblock}{ConfigurationMain.java}
		\begin{lstlisting}
SerializedSingleton singleton = SerializedSingleton.
    getInstance();
singleton.setValue(8);

try {
    FileOutputStream fileOut = new FileOutputStream("out.ser");
    ObjectOutputStream out = new ObjectOutputStream(fileOut);
    out.writeObject(singleton);
    out.close();
    fileOut.close();
} catch (IOException e) {
    e.printStackTrace();
}
			    
		\end{lstlisting}
 	\end{exampleblock}
\end{frame}

\note{
	Wir serialisieren unseren Singleton in eine Datei.
}

\begin{frame}[fragile]{Deserialisierung}
	\begin{exampleblock}{ConfigurationMain.java}
		\begin{lstlisting}
SerializedSingleton singleton2 = null;
try {
    FileInputStream fileIn = new FileInputStream("out.ser");
    ObjectInputStream in = new ObjectInputStream(fileIn);
    singleton2 = (SerializedSingleton) in.readObject();
    in.close();
    fileIn.close();
} catch (IOException i) {
    i.printStackTrace();
} catch (ClassNotFoundException c) {
    System.out.println("Singleton nicht gefunden");
    c.printStackTrace();
}
		\end{lstlisting}
 	\end{exampleblock}
\end{frame}

\note{
	Dann wird der Singleton deserialisiert.
}

\begin{frame}[fragile]{Vergleich der Objekte}
	\begin{alertblock}{ConfigurationMain.java}
		\begin{lstlisting}
if (singleton == singleton2) {
    System.out.println("Die beiden Objekte sind dieselben");
} else {
    System.out.println("Die beiden Objekte sind nicht dieselben");
}

singleton2.setValue(99);

System.out.println(singleton.getValue());
System.out.println(singleton2.getValue());
		\end{lstlisting}
 	\end{alertblock}
\end{frame}

\begin{frame}{Vergleich der Objekte}
  \begin{alertblock}{Ausgabe}
Die beiden Objekte sind nicht dieselben

8

99
  \end{alertblock}
\end{frame}

\note{
	Und tatsächlich erhalten wir zwei unterschiedliche Objekte trotz des privaten Konstruktors. Damit sind die Bedingungen für einen Singleton verletzt. Das Problem mit der Serialisierung im Singleton-Entwurfsmuster ergibt sich, weil die Standard-Serialisierung und Deserialisierung in Java ein Singleton-Objekt umgehen können, wodurch mehr als eine Instanz der Klasse entstehen kann. Beim Serialisieren wird der Zustand eines Objekts in einen Bytestrom umgewandelt, um ihn später zu speichern oder zu übertragen. Beim Deserialisieren wird aus diesem Bytestrom wieder ein Objekt erzeugt.
}

\begin{frame}[fragile]{Lösung}
	\begin{exampleblock}{SerializedSingleton.java}
		\begin{lstlisting}
public class SerializedSingleton implements Serializable {
	private int value;
	private static SerializedSingleton instance;
	
	private SerializedSingleton() {
	}

 	protected Object readResolve() {
		return instance;
	}
	
	public static SerializedSingleton getInstance() {
		if (instance == null) {
			instance = new SerializedSingleton();
		}
		
		return instance;
	}
}
		\end{lstlisting}
 	\end{exampleblock}
\end{frame}

\begin{frame}[fragile]{Lösung}
  \begin{exampleblock}{Ausgabe}
Die beiden Objekte sind dieselben

99

99
  \end{exampleblock}
\end{frame}

\note{
	Um das Problem zu lösen, muss die Methode readResolve() implementiert werden. Sie bereitet ein Objekt nach der Deserialisierung vor und stellt in unserem Fall die Singleton-Eigenschaft wieder her. 
}
